\begin{frame}{5항목의 구체적 의미 (3)-(5) 정당화 · 정직성 · 탐험-활용}
  \small
  \begin{itemize}
    \item \kb{(3) 수식적 정당화} — ``최소 1개''를 채웠는지는 \emph{형식}
      문제, 이 항목은 \emph{내용}: 짚은 개념(예: Ch11.3 GAE $\lambda$)이
      \textbf{실제로 관찰한 숫자와 일치}하는가. ``$\lambda$가 편향-분산을
      조절한다''는 문장만 있고 $\lambda$를 바꿔 \emph{어드밴티지 분산이
      어떻게 변했는지} 실측이 없으면, 그 문장은 \emph{어디에나 붙을}
      설명일 뿐.
    \item \kb{(4) 결과의 정직성} — 신경망 \textbf{초기화 시드}를 몇 개
      돌렸는지, 시드별 이동평균이 \emph{사라지지 않고} 보고에 들어갔는지,
      ``완전 수렴 아님''을 숨기지 않는지. ``결과가 좋았는가''가 아니라
      ``결과가 \textbf{정직한 절차}로 얻어졌는가''를 본다.
    \item \kb{(5) 탐험-활용 균형} — Ch02 멀티암 밴딧의 딜레마가
      \textbf{실제 학습 곡선}에서 어떻게 드러나는지. 신경망 RL에서는
      탐험이 $\varepsilon$-greedy(DQN, Ch09.2)나 가우시안 정책의 $\sigma$
      (Ch10.2, Ch11.2)로 구현 — ``탐험을 줄여가는데 줄이는 시점이 학습
      곡선의 어떤 지점이었나? 일찍 줄이면 고착, 늦으면 노이즈''(Ch02
      직관의 신경망 버전 검증).
  \end{itemize}
\end{frame}
